\chapter{Conclusions and Future Work}
\label{cap:conclusions}

This chapter summarizes the main findings, discusses the contributions and limitations of the work, and identifies directions for future research.

\section{Main contributions}
\label{sec:contributions}

This thesis investigated vision-based local guidance for agricultural row navigation through segmentation model selection, dataset preparation with model-assisted annotation, and geometric reference extraction. The main contributions are:

\begin{itemize}
\item \textbf{Architecture comparison for robotic segmentation:} FCN, DeepLabV3, and LR-ASPP were compared on manually annotated lettuce images. LR-ASPP was selected for its practical accuracy--speed trade-off, achieving path IoU of 0.9546 with 10.41~ms inference time compared to 58.24~ms and 86.87~ms for the other architectures. This 5--8$\times$ latency reduction supports real-time operation despite a 1.2 percentage point mIoU deficit.

\item \textbf{SAM~3-assisted dataset preparation:} Multi-crop training datasets for lettuce, tomato, and cucumber were prepared using SAM~3 text-prompted segmentation. This reduced manual annotation effort, although systematic measurements of time savings and inspection requirements remain to be completed.

\item \textbf{Evaluation of output label configurations:} Models with four-class, three-class, and binary path outputs were compared. The binary configuration aligns directly with the navigation task and was selected for robotic integration.

\item \textbf{Edge-based geometric extractors:} Two methods for converting corridor masks into image-space navigation references were implemented. Both use horizontal strip extraction and independent fitting of left and right boundaries, differing in seed selection, association strategy, and validation checks.
\end{itemize}

\section{Conclusions on segmentation and annotation}
\label{sec:segmentation_conclusions}

The architecture comparison establishes that LR-ASPP is suitable for real-time agricultural navigation when inference speed is a constraint. Its segmentation accuracy on lettuce is sufficient for path extraction, and its reduced computational cost makes it deployable on robot hardware.

SAM~3-assisted annotation demonstrates that foundation models can reduce manual labeling effort for agricultural datasets. The quality of automatically generated masks and their effect on trained model performance depend on crop appearance, text prompt formulation, and the need for manual inspection. Future work should measure annotation time systematically, including review and correction, to quantify the actual savings.

The choice of output label configuration affects both annotation effort and downstream geometric extraction. Binary path segmentation focuses the model directly on the navigation target, simplifying both the labeling task and the geometric stage. This comes at the cost of discarding plant and covering-material information that might be useful for other tasks.

\section{Conclusions on geometric extraction}
\label{sec:geometric_conclusions}

The geometric extraction stage converts binary corridor segmentation masks into navigation setpoints. The comparative study examined four distinct architectural paradigms:
\begin{enumerate}
    \item Global column projection (\texttt{pf\_histogram}) maximizes simplicity and eliminates tracking failure modes, but discards heading angle information and exhibits sensitivity to asymmetric vegetation distribution.
    \item Centerline strip regression (\texttt{pf\_centri}) provides both lateral offset and heading estimates with low computational demand, but transfers unilateral canopy noise directly into the navigation line.
    \item Random sample consensus (\texttt{pf\_ransac}) handles scattered boundary points and wide gaps without strip discretization, but introduces frame-to-frame stochastic jitter ($\sigma_\omega > 0$) and variable runtime.
    \item Decoupled dual-edge tracking (Method~4) provides deterministic execution ($\sigma_\omega = 0.00^\circ$), sub-millisecond latency, and boundary-level noise decoupling, but introduces strict dependencies on near-field seed detection and heuristic jump gates.
\end{enumerate}

These characteristics demonstrate that no individual extractor is universally superior; selection depends on sensor setup, platform compute limits, and field layout.

\todobox{Sintesi quantitativa finale dell'estrazione geometrica e navigazione}{%
Sintetizzare i dati conclusivi:
\begin{itemize}
  \item Margine di precisione medio dell'asse stimato (scostamento in pixel e in centimetri).
  \item Tasso medio di affidabilit\`a e di reiezione dei frame lungo i test eseguiti.
  \item Conferma della sostenibilit\`a del controllo a 20 Hz su piattaforma embedded Jetson in condizioni di campo.
\end{itemize}
}

\section{Limitations}
\label{sec:limitations}

The methodology and experimental findings are subject to several specific limitations:

\begin{itemize}
\item \textbf{Bottom-up seed dependency in Method~4:} The tracking algorithm initializes exclusively from the bottom six horizontal strips. Occlusions immediately in front of the camera (such as robot chassis shadows, soil buildup, or puddles) prevent track initialization and lead to frame rejections, even when the distant corridor is visible.

\item \textbf{Geometric model constraints:} Method~4, \texttt{pf\_ransac}, and \texttt{pf\_centri} assume locally straight, parallel corridor boundaries. They cannot follow curved crop rows or navigate headland turnaround maneuvers where crop boundaries terminate or diverge.

\item \textbf{Unilateral row dropout vulnerability:} If crop continuity along one row is broken by gaps exceeding three strips ($18$~px), Method~4 drops that boundary. Without an asymmetric single-edge fallback mode, this results in reference rejection, whereas global histogram projection continues to provide a lateral setpoint.

\item \textbf{Image-space gate tuning:} The continuity gate ($\Delta x_{\text{max}} = 28$~px) is defined in image coordinates. Significant vehicle pitch and roll on uneven soil cause sudden lateral displacements of the corridor between consecutive scanlines, which can trigger premature tracking loss.

\item \textbf{Dataset annotation coverage:} Segmentation models were quantitatively evaluated against dense manual ground truth on lettuce. Multi-crop training includes tomato and cucumber, but their segmentation accuracy was not verified against dense manual annotations for those specific crops.

\item \textbf{Annotation effort measurement:} Time savings from SAM~3-assisted annotation were evaluated qualitatively. The comparison does not report logged inspection hours or the frequency of manual polygon corrections.

\item \textbf{Closed-loop dynamics:} Offline benchmarks evaluate reference accuracy on static test images. They do not account for vehicle inertia, steering backlash, or dynamic latency variations under closed-loop field control.
\end{itemize}

\section{Future work}
\label{sec:future_work}

Several directions for future research emerge from this work:

\subsection{Dataset and annotation}
\label{subsec:future_annotation}

\begin{itemize}
\item Measure annotation time systematically for manual and SAM~3-assisted workflows, including inspection and correction.
\item Prepare manual ground-truth masks for tomato and cucumber to validate multi-crop model performance directly.
\item Investigate active learning strategies to select images for manual annotation that improve model generalization most effectively.
\item Explore other foundation models or fine-tuned SAM variants specialized for agricultural segmentation.
\end{itemize}

\subsection{Segmentation and model deployment}
\label{subsec:future_segmentation}

\begin{itemize}
\item Evaluate additional lightweight architectures (e.g., BiSeNet, FastSCNN) for further latency reduction.
\item Investigate quantization and model compression techniques to enable deployment on edge devices.
\item Test temporal consistency methods that use previous frames to improve segmentation and reduce flickering.
\item Explore uncertainty estimation to detect when segmentation confidence is low and fallback strategies are needed.
\end{itemize}

\subsection{Geometric extraction}
\label{subsec:future_geometric}

\begin{itemize}
\item Extend the geometric representation from linear models to piecewise linear or polynomial formulations to support curved rows and headland turns.
\item Implement top-down recovery logic and single-edge asymmetric tracking to maintain guidance when bottom-up seeding fails or unilateral row dropouts occur.
\item Adapt spatial gate thresholds dynamically using vehicle attitude telemetry from an onboard IMU to prevent false tracking loss during pitch and roll variations.
\item Incorporate temporal Bayesian tracking (such as an Extended Kalman Filter) to smooth reference parameters across frame dropouts without incurring the stochastic jitter of RANSAC.
\end{itemize}

\subsection{System integration and field validation}
\label{subsec:future_integration}

\begin{itemize}
\item Complete closed-loop field trials to validate end-to-end navigation performance under variable lighting, growth stages, and row conditions.
\item Measure lateral tracking error and success rate over extended distances and multiple crop types.
\item Integrate temporal filtering and prediction to handle brief reference loss or occlusions.
\item Investigate fusion with other sensors (e.g., depth, IMU) to improve robustness.
\item Develop a complete autonomy system including global path planning, obstacle avoidance, and task-specific behaviors.
\end{itemize}

\section{Closing remarks}
\label{sec:closing}

This thesis addressed the segmentation, dataset preparation, and geometric extraction stages of vision-based agricultural navigation. The selection of LR-ASPP balances accuracy and speed for robotic deployment. SAM~3-assisted annotation reduces manual effort, and the geometric extractors convert masks into navigation references using horizontal strip selection and independent edge fitting.

The work establishes a foundation for camera-based local guidance in structured crop environments. Future research should complete field validation, extend the geometric model to handle curves and gaps, and systematically measure annotation savings to guide dataset preparation strategies.
